\scanpage{15}
Pour tte section $\unsure{\zeta|\xi}$ de $\check E$ (i.e hom $\varphi:M\to\mathbb G_a$) on peut considérer les \unsure{sections|section}
\[
 \zeta^{(p)\unsure{*|1}}\overset{\mathrm{def}}{=}
 u(\zeta^{(F)\,\to\,\text{frobénius})})
\]
et par itération les $\zeta^{(p^i)}$. Elles correspondent aux composés de
$\varphi:M\mathrel{\rule{1.3em}{.4pt}}G_a$ avec $F_{G_a}:G_a\to G_a$,
et \deleted{\unsure{leurs|les}} \inserted{ses} itérés $F_{G_a}^{i}$.

Considérons
\[
 \zeta\wedge\zeta^{(p)}\ldots\wedge\zeta^{(p^{n-1})}
 \in\Gamma(\textstyle\bigwedge^n\check E)
\]
On trouve ainsi \inserted{(en changeant de base)} une application \unsure{polynômiale|polynômial|polynomiale}
\[
 \Delta^{E}_{\unsure{M|n}}:\check E\longrightarrow\textstyle\bigwedge^n\check E
\]
homogène de degré $1+p+\text{---}+p^{\unsure{n-1}}=\frac{p^n-1}{p-1}$, d'où
\deleted{d'où \unsure{un morphisme}}
\[
 \begin{aligned}
 \Delta_{n}^E:\quad &\mathrm V(E)\longrightarrow \mathrm V(\textstyle\bigwedge^{n}E)\\
 &\hspace{.7em}\shortparallel \hspace{6em}\shortparallel\\[-2pt]
 &\mathrm W(\check E)\hspace{3.2em}\mathrm W(\textstyle\bigwedge^{n}\check E)
 \end{aligned}
\]
\deleted{\uline{Prop} \unsure{Sur} chaque fibre,}

Considérons l'image inverse dans $\mathrm V(E)$ de la section \unsure{nulles|nulle} de
$\mathrm V(\bigwedge^{n}E)$, \deleted{\unsure{soit} \illegible{1}}
c'est un diviseur de Cartier relatif, dont le complémentaire
$\mathrm V(\deleted{\unsure{\bigwedge}}E)^0$ est un ouvert \unsure{univt} sch. dense.
Si $\zeta$ est une section de cet ouvert, alors $\zeta,\zeta^{(p)}\ldots\zeta^{(p^{n-1})}$
\deleted{\unsure{forment}} une base de $\check E$, et on peut écrire
\[
 \zeta^{(p^n)}+A^E_{n-1}(\zeta)\,\zeta^{(p^{n-1})}
 \unsure{+\cdots+|-\cdots+}A^E_0(\zeta)\,\zeta=0
\]
% archive sheet number: 14
